% ===================================================================
% ABSTRACT (MACEDONIAN)
% ===================================================================
\chapter*{Апстракт}
\addcontentsline{toc}{chapter}{Апстракт}

\textbf{Вовед.} Федерираното учење овозможува системите за откривање упади во мрежите на
Интернетот на нештата да се обучуваат без пренос на суров сообраќај, но го отвора прашањето
на нивното одржување: кога еден јазол ќе биде идентификуван како компромитиран, или кога
правото на заборавање ќе биде искористено, влијанието на неговите податоци мора да биде
отстрането од моделот. Бришењето на суровите записи не е доволно, бидејќи моделот ги
задржува во своите тежини информациите научени од нив. Постојната литература за федерирано
заборавање ја изразува цената на отстранувањето во епохи и рунди, и ја мери во симулација.
На вистински уред на работ на мрежата --- ограничен истовремено по пресметка, топлина,
складирање и напојување --- таа апстракција не е доволна. Овој труд ја мери \textbf{физичката цена
на заборавањето} на реален хардвер.

\textbf{Методи.} Изградена е хибридна федерација од четири клиенти врз рамката Flower 1.29 во
режим на распоредување, при што централниот агрегатор и три симулирани јазли работат на
домаќин, а четвртиот јазол е физички Raspberry Pi 5 (4 GB) приклучен преку локална мрежа.
Уредот е намерно ограничен: вентилаторот е исклучен, што наметнува термичко придушување, а
складирањето е SD-картичка од класата A1. Потрошувачката се мери со надворешен мерач FNIRSI
FNB58 поставен сериски на напојувањето. Врз податочното множество TON\_IoT, со по 1\,000\,000
реда по јазол, се споредуваат два пристапа за егзактно локално заборавање: \textbf{наивно повторно
тренирање} од нула врз задржаните податоци и \textbf{заборавање со SISA} ($S = 5$ шардови
× $R = 5$ сегменти) со враќање наназад и повторна обработка само на засегнатите сегменти.
Отстранувањето е предизвикано од труење со превртување на етикети, ограничено на сегментите
3–4 на еден шард, кое дефинира јасно определено множество од околу 63\,000 примери што мора
да се заборават. Дизајнот е претходно регистриран и замрзнат пред мерењата, спарен по
почетна вредност, со $N = 10$ пара (42–51) и балансиран редослед, а статистиката е
непараметриска: Вилкоксонов тест на рангирани знаци, Cliff's $\delta$ и bootstrap 95\,\%
интервали на доверба.

\textbf{Резултати.} SISA го намалува времето до обновување од медијана 242,1 s на 29,5 s
(\textbf{8,2×}), нето-потрошената енергија од 0,093 Wh на 0,017 Wh (\textbf{5,4×}) и времето поминато
во термичко придушување од 233 s на 27,5 s (\textbf{8,5×}). Сите три достигнуваат
$p = 0{,}0020$ --- долната граница при $N = 10$ --- со Cliff's $\delta = +1{,}00$, односно
целосно раздвојување: ниту едно извршување со SISA не се преклопува со ниту едно наивно
извршување. Трошокот се поместува кон складирањето (12,05 наспроти 2,75 MB запишани), но
апсолутната големина е занемарлива за модел со 12\,865 параметри. Режискиот трошок на SISA во
фазата на обука изнесува 43,67 MB и медијана од 5,6 s, што едно единствено обновување го
отплаќа повеќе од триесеткратно. \textbf{Корисноста бара внимателно читање.} Повторната евалуација
на секој зачуван глобален модел врз издвоено тест-множество од 100\,000 примери со 96,5\,\%
напади покажува дека при подразбираниот праг од 0,5 обновениот модел со SISA предвидува дека
секој тек е напад: специфичност 0,000, избалансирана точност 0,500 (случаен избор) и MCC
0,017, наспроти 0,237, 0,598 и 0,177 за наивниот пристап; одѕивот од 1,000 и $F_1$ од 0,982
се тривијални последици на класната нерамнотежа. Колапсот, меѓутоа, е \textbf{калибрациски}: со
подесен праг SISA достигнува избалансирана точност \textbf{0,709 наспроти 0,671} за наивниот
пристап, подобра во 9 од 10 почетни вредности и истовремено по специфичност и по MCC. Тој
му претходи на обновувањето и се локализира во параметарското просечување на составните
модели --- поединечно тие дискриминираат, средната вредност на нивните тежини не.

\textbf{Заклучок.} SISA го прави заборавањето на уред на работ осумкратно поевтино, со занемарлив
трошок на подготовка и \textbf{без загуба во корисноста --- под услов прагот на одлучување да биде
калибриран}. Некалибрирана, таа испорачува детектор што сè прогласува за напад: дефект што
стандардните показатели одѕив и $F_1$ целосно го скриваат при изразена класна нерамнотежа, и
кој произлегува од агрегацијата, не од заборавањето. Придонесот на трудот е методологија за
мерење на физичката цена на заборавањето и емпириска основа за избор меѓу двата пристапа,
изразена во величини што проектантот на системи навистина ги плаќа.

\textbf{Клучни зборови:} федерирано учење, машинско заборавање, SISA, откривање упади, Интернет
на нештата, пресметување на работ, енергетска потрошувачка, термичко придушување, класна
нерамнотежа, право на заборавање.

